% Generated draft: requires mathematical and source review.
\documentclass[11pt]{article}
\usepackage[margin=25mm]{geometry}
\usepackage{fontspec}
\setmainfont{Latin Modern Roman}
\usepackage{amsmath,amssymb,mathtools,unicode-math}
\setmathfont{Latin Modern Math}
\usepackage{enumitem,needspace,xurl,hyperref}
\setlength{\parindent}{0pt}
\setlength{\parskip}{5pt}
\setlength{\emergencystretch}{3em}
\newcommand{\N}{\mathbb{N}}
\newcommand{\Z}{\mathbb{Z}}
\newcommand{\Q}{\mathbb{Q}}
\newcommand{\R}{\mathbb{R}}
\newcommand{\C}{\mathbb{C}}
\newcommand{\F}{\mathbb{F}}
\newcommand{\A}{\mathbb{A}}
\newcommand{\PP}{\mathbb{P}}
\newcommand{\E}{\mathbb{E}}
\newcommand{\eps}{\varepsilon}
\newcommand{\dd}{\,\mathrm{d}}
\begin{document}
\section*{UnsolvedMath ID: 30001251}
\textbf{Source problem code:} catalogue entry
\subsection{Definitions and mathematical statement}
Let \(S^1=\{z\in\C:|z|=1\}\). A steady state means an admissible solution of the evolution equation that is independent of time. The supplied record refers to a noisy orientational aggregation model on \(S^1\), but does not specify its equation, admissible states, interaction law, noise mechanism, or hypotheses ensuring the existence of states with \(N\) peaks. It also does not define a peak or the intended notion of instability.

\textbf{Conjecture as supplied.} For an integer \(N\geq 3\), under the hypotheses ensuring the existence of \(N\)-peaked steady states of the noisy orientational aggregation model on \(S^1\) for sufficiently small noise, are all such steady states unstable?

For a schematic expression only, let \(\eta\) denote the unspecified noise parameter and let \(\mathcal S_N(\eta)\) denote the family of steady states described in the source as \(N\)-peaked, subject to its unspecified existence hypotheses. The proposed assertion is
\[
\forall u_*\in\mathcal S_N(\eta),\qquad u_*\text{ is unstable},
\]
for each integer \(N\geq3\) and sufficiently small noise under those hypotheses. These symbols do not supply the missing model or stability definition. The record does not determine whether the small-noise bound may depend on \(N\), the model parameters, or the particular steady-state family. Consequently, a fully specified mathematical conjecture cannot be recovered from this record alone.
\subsection{Short English statement}
The supplied conjecture asks whether every steady state with at least three peaks is unstable in the small-noise regime where the stated existence hypotheses apply. The model, existence hypotheses, meaning of a peak, and stability criterion are missing, so this remains a qualified statement requiring source consultation and human mathematical review.
\subsection{Sources}
\begin{itemize}
\item Supplied local UnsolvedMath registry record: permanent ID \(30001251\), title \emph{Instability of Multi-Peaked Orientational Steady States}, source problem code \texttt{catalogue entry}, in the Oberwolfach Reports Open Problems collection.
\item Supplied report citation: \emph{Mathematical Biology} (2009), Oberwolfach Reports, pp.\ 1303--1374, DOI \texttt{10.4171/owr/2009/24}. Supplied source URL and literature-source link, labelled ``original Oberwolfach report'': \url{https://doi.org/10.4171/owr/2009/24}. The background attributes the collection to mathematical workshop reports and their open-problem discussions; no individual proposer is supplied.
\item Literature assessment dated \texttt{2026-08-21}: the supplied curation labels the problem open and describes the assertion of instability for all such \(N\)-peaked states as a conjecture supported only by numerical evidence; it reports that no proof or counterexample was identified. This is supplied dated curation, not an independently verified literature assessment. The supplied background requests independent verification of time-sensitive status and cited claims.
\item The registry labels the statement \texttt{corrected\_verified}, explaining that this denotes meaning-preserving semantic, typographic, and LaTeX cleanup without unsupported source-dependent additions. This supplied label does not establish mathematical correctness or independent verification of the conjecture.
\end{itemize}
\end{document}
